A afirmação de que são necessários nove gráficos, cada um contendo três curvas,
para relacionar força e velocidade, pode ser justificada da seguinte forma.

\begin{enumerate}
    \item \textbf{Número de parâmetros variáveis.}
    Suponha que a força dependa de três parâmetros independentes (por exemplo,
    massa, densidade e área). Se cada parâmetro puder assumir três níveis
    distintos (baixo, médio e alto), então o número total de combinações
    possíveis é
    \[
        3^3 = 27.
    \]

    \item \textbf{Organização das combinações.}
    Para representar graficamente essas 27 combinações de forma clara, agrupa-se
    cada conjunto de três combinações que diferem apenas em um dos parâmetros.
    Assim, cada gráfico contém três curvas, correspondentes às três variações
    possíveis de um único parâmetro, enquanto os outros dois permanecem fixos.

    \item \textbf{Número total de gráficos.}
    Como existem 27 combinações e cada gráfico comporta três curvas, o número
    total de gráficos necessários é
    \[
        \frac{27}{3} = 9.
    \]

    \item \textbf{Clareza na análise.}
    Essa organização permite comparar diretamente o efeito de um único
    parâmetro sobre a relação força–velocidade, mantendo os demais constantes.
    Dessa forma, a interpretação física torna-se mais simples e visualmente
    clara.
\end{enumerate}

Portanto, são necessários nove gráficos, cada um contendo três curvas, para
representar todas as combinações possíveis dos três parâmetros e analisar
adequadamente a relação entre força e velocidade.